
% ════════════════════════════════════════════
% 具体分析（5.2.1.具体分析.tex，三级子块）
% ════════════════════════════════════════════
\subsubsection{不确定性的来源与三段式框架}

针对问题二，本问属于\textbf{在预测不确定性下、含两阶段决策的全年逐日滚动优化}问题，与问题一存在两点本质区别：题目要求在每天 0:00 制定当天的计划购电量，而当天 144 个间隔的负载与光伏功率在 0:00 时尚未发生，若在优化中直接引用附件二给出的实际功率，则相当于使用了事后信息，模型会认为供电缺口恒为零、紧急购电永远用不到，这与题目专门设置 5 倍紧急电价并要求以表 3 格式给出紧急购电时段与电量的意图相矛盾；同时，预测必然存在误差，误差会以供电缺口的形式暴露出来，只能依靠 5 倍电价的紧急购电兜底。因此“预测”是本问绕不开的独立子问题，必须显式建模。

基于上述判断，本文按\textbf{“先预测、再决策、后结算”三段式}组织求解：\textbf{预测段}以 1 月实际数据为训练样本，逐日扩张窗口标定负载与光伏预测模型，在每天 0:00 给出当天 144 个间隔的功率预测值，并由历史预测残差估计其不确定性；\textbf{决策段}按负载低/中/高与光伏低/中/高组合构造 9 个情景，建立两阶段\textbf{随机线性规划}，一阶段决策（计划购电量与储能充放电计划）要求对所有情景同时可行，故与当日真值无关，目标函数为计划购电费与 5 倍电价下的期望紧急购电费之和；\textbf{结算段}在真实功率到达后，按既定充放电计划计算实际供电缺口，以 5 倍电价结算紧急购电量，富余电量计为弃电，日末储电量传递给次日。计算以 2025 年 1 月 1 日 0:00 储电量 6000 kWh 为起点逐日滚动，全年 365 天中 1 月作为预热期，最终输出 2025 年 2 月 1 日至 12 月 31 日共 334 天的计划购电量、储能充放电量、紧急购电量与弃电量。具体分析流程如图\ref{fig:analysis2_flow}所示。

\begin{figure}[H]
    \begin{center}
        \begin{tikzpicture}[x=1cm, y=0.75cm,
            flowhead/.style={rectangle, rounded corners=3pt, draw=green!60!black, fill=green!45!black, text=white, font=\bfseries, minimum width=8.4cm, minimum height=1.1cm, align=center, inner sep=2pt, line width=0.7pt},
            flowbox/.style={rectangle, rounded corners=2.5pt, draw=green!55!black, fill=green!12, minimum width=8.4cm, minimum height=1.0cm, align=center, inner sep=2pt, line width=0.6pt},
            flowend/.style={rectangle, rounded corners=2.5pt, draw=green!65!black, fill=green!35!black, text=white, font=\bfseries, minimum width=8.4cm, minimum height=1.0cm, align=center, inner sep=2pt, line width=0.7pt},
            flowarrow/.style={thick, -{Stealth}, green!50!black}
        ]
            \node[flowhead] at (0, 0) (h) {问题二：预测—决策—结算三段式};
            \node[flowbox] at (0, -1.8) (s1) {\textbf{读取全年数据}\\\footnotesize 365 天负荷 · 光伏 · 分时电价};
            \node[flowbox] at (0, -3.6) (s2) {\textbf{标定预测模型}\\\footnotesize 1 月扩张窗口 · 形态 + 水平 + 残差池};
            \node[flowbox] at (0, -5.4) (s3) {\textbf{0:00 预测当日功率}\\\footnotesize 144 间隔负载与光伏预测值};
            \node[flowbox] at (0, -7.2) (s4) {\textbf{9 情景两阶段随机规划}\\\footnotesize 计划购电与充放电计划·与真值无关};
            \node[flowbox] at (0, -9.0) (s5) {\textbf{按计划执行并结算}\\\footnotesize 真实功率到达 · 计算供电缺口};
            \node[flowbox] at (0, -10.8) (s6) {\textbf{5 倍电价紧急购电}\\\footnotesize 缺口计紧急购电 · 富余计为弃电};
            \node[flowbox] at (0, -12.6) (s7) {\textbf{日末储电量传递}\\\footnotesize 当日末储电量作为次日初值};
            \node[flowend] at (0, -14.4) (s8) {汇总输出结果\\\footnotesize 2 月 1 日—12 月 31 日计划/紧急/总费用};
            \draw[flowarrow] (h) -- (s1);
            \draw[flowarrow] (s1) -- (s2);
            \draw[flowarrow] (s2) -- (s3);
            \draw[flowarrow] (s3) -- (s4);
            \draw[flowarrow] (s4) -- (s5);
            \draw[flowarrow] (s5) -- (s6);
            \draw[flowarrow] (s6) -- (s7);
            \draw[flowarrow] (s7) -- (s8);
        \end{tikzpicture}
        \caption{问题二分析流程图}
        \label{fig:analysis2_flow}
    \end{center}
\end{figure}
